\documentclass[a4paper,12pt,reqno]{amsart}
\usepackage{amsfonts}
\usepackage{amsmath}
\usepackage{amssymb}
\usepackage[a4paper]{geometry}
\usepackage{mathrsfs}
\usepackage{csquotes}
\usepackage{enumitem}

%\usepackage{esint}
%
%\usepackage{draftcopy}
\usepackage{xcolor}
\usepackage[colorlinks = true,
            linkcolor = red,
            urlcolor  = gray,
            citecolor = blue,
            anchorcolor = blue]{hyperref}
\renewcommand\eqref[1]{(\ref{#1})} %Need with hyperref
%
%
%%%%%%%%%%%%%%%%%%%%%%%%%
\setlength{\textwidth}{15.2cm}
\setlength{\textheight}{22.7cm}
\setlength{\topmargin}{0mm}
\setlength{\oddsidemargin}{3mm}
\setlength{\evensidemargin}{3mm}
%\setlength{\headsep}{2.9cm}
\setlength{\footskip}{1cm}

%%%%%%%%%%%%%%%%%%%%%%%%%%%

\numberwithin{equation}{section}
\theoremstyle{plain}
\newtheorem{thm}{Theorem}[section]
\newtheorem{prop}[thm]{Proposition}
\newtheorem{cor}[thm]{Corollary}
\newtheorem{lem}[thm]{Lemma}
 \newtheorem{exa}[thm]{Example}
\theoremstyle{definition}
\newtheorem{defn}[thm]{Definition}
\newtheorem{rem}[thm]{Remark}
\newtheorem{ex}[thm]{Example}
\newtheorem{property}[thm]{Property}
\renewcommand{\wp}{\mathfrak S}
\newcommand{\Rn}{\mathbb R^{n}}

\def\G{{\mathbb G}}
\def\R{{\mathcal R}}
\def\L{{\mathcal L}}
\DeclareMathOperator*{\esssup}{ess\,sup}
%%%%%%%%%%%%%%%%%%%%%%%%%%%%%

%   \documentclass{rmmcart}

\begin{document}


   \title[Nonlinear spectral problem  for H\"ormander vector fields]
   {Nonlinear spectral problem \\ for H\"ormander vector fields}

\author[M. Karazym]{Mukhtar Karazym}
\address{
  Mukhtar Karazym:
  \endgraf
  Department of Mathematics
  \endgraf
Nazarbayev University, Kazakhstan
  \endgraf
  {\it E-mail address} {\rm mukhtarkarazym@gmail.com; mukhtar.karazym@nu.edu.kz}
  }

 \author[D. Suragan]{Durvudkhan Suragan}
\address{
	Durvudkhan Suragan:
	\endgraf
	Department of Mathematics
	\endgraf
Nazarbayev University, Kazakhstan
	\endgraf
	{\it E-mail address} {\rm durvudkhan.suragan@nu.edu.kz}
} 



\keywords{nonlinear spectral problem, variational method, H\"ormander vector field,  Carnot-Carath\'eodory space}
\subjclass[2020]{35P30, 35A15, 53C17}

\begin{abstract} 
Based on variational methods, we study a nonlinear eigenvalue problem for a quasilinear operator arising from smooth H\"ormander vector fields. We derive the smallest eigenvalue, prove its simplicity and  isolatedness, establish  the positivity of the first eigenfunction and show  H\"older regularity of eigenfunctions.  Moreover, we determine the best constant for the $L^{p}$-Poincar\'e inequality as a byproduct.
\end{abstract}
\maketitle

%\tableofcontents

\section{Introduction}
The following nonlinear eigenvalue problem represents the simplest prototype of our objective:
\begin{equation}\label{Dirichlet p-Laplacian}
\begin{aligned}
\operatorname{div}\left(|\nabla u|^{p-2} \nabla u\right)&=-\lambda|u|^{p-2} u,  &&\text { in } \Omega, \\
u&=0,
&&\text { on } \partial \Omega,
\end{aligned}
\end{equation}
where $1<p<\infty$, $\nabla$ is the usual gradient and $\Omega\subset\mathbb{R}^{n}$ is a bounded open set with smooth boundary $\partial\Omega$.  Extensive literature exists on problem \eqref{Dirichlet p-Laplacian}. However, relatively limited research has been conducted on nonlinear eigenvalue problems for the $p$-sub-Laplacian on a stratified Lie group $\mathbb{G}=\left(\mathbb{R}^{n}, \circ\right)$ (with the group operation $\circ$)
\begin{equation}\label{Dirichlet p-sub-Laplacian}
\begin{aligned}
\nabla_{\mathbb{G}} \cdot\left(\left|\nabla_{\mathbb{G}} u\right|^{p-2} \nabla_{\mathbb{G}} u\right)&=-\lambda|u|^{p-2} u,  &&\text { in } \Omega\subset \mathbb{G}, \\
u&=0,
&&\text { on } \partial \Omega,
\end{aligned}
\end{equation}
where $\nabla_{\mathbb{G}}$ is the horizontal gradient on $\mathbb{G}$.

In a previous study \cite{WNL2009}, several estimates of the first and second eigenvalues for the Dirichlet $p$-sub-Laplacian \eqref{Dirichlet p-sub-Laplacian} were obtained. Recently, Carfagnini and Gordina \cite{CG2022} focused on the linear case of the spectral problem \eqref{Dirichlet p-sub-Laplacian} when $p=2$, thereby advancing our understanding of this topic within stratified Lie groups.

Chen, et.al. \cite{Chen2021}, assuming only the H\"ormander finite rank condition for smooth vector fields on $\Omega$, studied the eigenvalue problem in Carnot-Carath\'eodory spaces:
\begin{equation}\label{Dirichlet sub-Laplacian type}
\begin{aligned}
\Delta_{X} u&=-\lambda u,  &&\text { in } \Omega, \\
u&=0,
&&\text { on } \partial \Omega.
\end{aligned}
\end{equation}
It is worth mentioning that, when $p=2$, the authors found the best constant of the Poincar\'e inequality in \cite{Chen2021}. In this paper, we extend the finding of Chen, et al. \cite{Chen2021} and establish the best constant for the entire range $1<p<\infty$ (see Remark \ref{the best constant of Poincare}).


Thus, the present paper aims to study a nonlinear eigenvalue problem for $p$-sub-Laplacian type operators in Carnot-Carath\'eodory spaces:
\begin{equation}\label{Dirichlet p-sub-Laplacian type}
\begin{aligned}
\sum_{j=1}^{m}X_{j}^{*}\left(\left|X u\right|^{p-2} X_{j}u\right)&=-\lambda|u|^{p-2} u,  &&\text { in } \Omega, \\
u&=0,
&&\text { on } \partial \Omega,
\end{aligned}
\end{equation}
where $X_{j}$ is a smooth  H\"ormander vector field defined over $\Omega$ and $X_{j}^{*}$ is the formal adjoint of $X_{j}$. The findings presented in this paper contribute to the understanding and analysis of eigenvalue problems associated with quasilinear operators \eqref{Dirichlet p-sub-Laplacian type} and provide insights into the behavior of eigenfunctions.

Our paper is organized as follows. Section \ref{sec2} provides a concise overview of H\"ormander vector fields, weak Sobolev spaces, and Sobolev embeddings. In Section \ref{sec3}, we delve into the study of a nonlinear eigenvalue problem concerning smooth H\"ormander vector fields. Our main results encompass the existence of the first eigenvalue and the positiveness of the corresponding eigenfunction. The positiveness result is derived from the application of the Harnack inequality. Similarly to the classical case, obtaining the first eigenvalue involves minimizing the Rayleigh quotient, leading to the determination of the best constant of the $L^{p}$-Poincar\'e inequality. Section \ref{sec4} focuses on investigating the simplicity and isolated nature of the first eigenvalue. The isolatedness property relies on the lower bound of eigenvalues (Lemma \ref{lower bound of eigenvalues}) and compact Sobolev embeddings (Theorem \ref{compact Sobolev embedding}).

\section{Preliminaries}\label{sec2}
Let $\Omega$ be a bounded connected open set in $\mathbb{R}^{n}$ with smooth boundary $\partial\Omega$. Consider a family of smooth vector fields $X_{1}, \ldots, X_{m}$  satisfying H\"ormander finite rank condition
\begin{equation*}
\operatorname{rank}\left(\operatorname{Lie}\left\{X_{1}, \ldots, X_{m}\right\}\right)=n.
\end{equation*}
\begin{defn}
Given two points $x,y\in \Rn$,
the control distance $d_X(x, y)$ associated with the family of vector fields  $\{X_j\}$ is defined as the infimum of $T>0$ such that there is a piecewise $C^{1}$   curve $\gamma$ of $X_1, \ldots, X_m$ with values $\gamma(0)=x$ and $\gamma(T)=y$.   
\end{defn}
\begin{property}\label{Property 1}
Carnot-Carath\'eodory space $(\Omega, d_X)$ is complete and length-space by \cite[Lemma 3.1]{Str1986}.   
\end{property}

\begin{prop}\cite[p. 111]{NSW1985}\label{doubling condition}
For all $r>0$ and $x\in \Omega^{'}\Subset \Omega$ there exists a constant $C>0$ such that the doubling condition 
\begin{equation}\label{doubling}
    |B(x,2r)|\leq C |B(x,r)|
\end{equation}
holds, where
\begin{equation*}
B(x, r)=\{y \in \Omega \mid d_{X}(x, y)<r\} .
\end{equation*}
\end{prop}
\begin{defn}
The number 
\begin{equation}\label{homogeneous dimension Q}
    Q=\log_{2}C
\end{equation}
is called the local homogeneous dimension of $\Omega$, where $C$ is given in \eqref{doubling}.    
\end{defn}

The notation $X:=\left(X_{1}, \ldots, X_{m}\right)$ stands for the horizontal gradient
and $$|X u|=\left(\sum_{j=1}^{m}\left(X_{j} u\right)^{2}\right)^{\frac{1}{2}}.$$

We also recall a weak form of Jerison's Poincar\'e inequality \cite{Jer1986} 
\begin{equation}\label{L1 Poincare}
\int_B\left|u-u_B\right|^p d x \leq C r^p \int_{\alpha B}|X u|^p d x, \quad 1 \leq p<\infty ,
\end{equation}
where $B:=B(x,r)$, $\alpha B:=B\left(x, \alpha r\right)$ with $\alpha,r>0$ and $u_B=\frac{1}{|B|} \int_B u d x$ (see also \cite[p. 1122]{GN1996}).



The distributional derivative of a function $f$ along  $X_{j}$ is given by the identity 
\begin{equation*}
\left\langle X_{j} f, \varphi\right\rangle=\int_{\Omega} f X_{j}^{*} \varphi d x\quad \text{ for all } \varphi \in C_{c}^{\infty}(\Omega).
\end{equation*}

Throughout the paper we assume that $p>1$ unless otherwise stated.
\begin{defn}
We define a weak Sobolev space \begin{equation*}
\mathscr{L}^{1, p}(\Omega):=\left\{f \in L^{p}(\Omega): X_{j} f \in L^{p}(\Omega)\quad \text{for } j=1, \ldots, m \right\}
\end{equation*}
endowed with the norm
\begin{equation*}\label{natural norm}
\|u\|_{\mathscr{L}^{1,p}(\Omega)}:=\left(\int_{\Omega}\left|u|^{p}+|X u\right|^{p} d x\right)^{\frac{1}{p}}.
\end{equation*}
\end{defn}

\begin{defn}
The Sobolev space $S^{1,p}(\Omega)$ is the completion of $C^{\infty}(\Omega)$ with respect to the norm $\|\cdot\|_{\mathscr{L}^{1,p}(\Omega)}$.
\end{defn}
Note that $S^{1,p}(\Omega)=\mathscr{L}^{1, p}(\Omega)$ by \cite[Theorem 1.13]{GN1996}.
\begin{defn}
The Sobolev space $S_{0}^{1,p}(\Omega)$ is the completion of $C_{c}^{\infty}(\Omega)$ in $S^{1,p}(\Omega)$.
\end{defn}
\begin{rem}(Poincar\'e inequality). There exists a  constant $C>0$ such that the inequality
\begin{equation}\label{Poincare inequality}
\int_{\Omega}|u|^{p} d x\leq C \int_{\Omega}|X u|^{p} d x
\end{equation}
holds for all $u\in S_{0}^{1,p}(\Omega)$. Here the constant $C$ is independent of $u$.
\end{rem}

The space $S_{0}^{1,p}(\Omega)$ is a reflexive Banach space and can be equipped with the equivalent norm
\begin{equation}\label{Poincare norm}
\|u\|_{S_{0}^{1,p}(\Omega)}:=\left(\int_{\Omega}\left|X u\right|^{p} d x\right)^{\frac{1}{p}}.
\end{equation}

Property \ref{Property 1}, Proposition \ref{doubling condition} and Poincar\'e inequality \eqref{L1 Poincare}  satisfy hypotheses  in \cite[pp. 1084-1085]{GN1996}.
Therefore, the following compact Sobolev embeddings are valid.
\begin{thm}\label{compact Sobolev embedding}\cite[Theorem 1.28]{GN1996}
\begin{enumerate}
    \item For all $1\leq p <Q$ the embedding $S^{1,p}_{0}(\Omega)\hookrightarrow L^{q}(\Omega)$ is compact, where $1\leq q < \frac{Qp}{Q-p}$ and $Q$ is given by \eqref{homogeneous dimension Q}.  
    \item For all $Q\leq p  <\infty$ and $1\leq q <\infty$ the embedding $S^{1,p}_{0}(\Omega)\hookrightarrow L^{q}(\Omega)$ is compact.
\end{enumerate}\end{thm}

\section{Nonlinear eigenvalue problem}\label{sec3}
Consider the nonlinear eigenvalue problem
\begin{equation}\label{eigenvalue problem}
\begin{aligned}
X \cdot\left(\left|X u\right|^{p-2} X u\right)&=-\lambda|u|^{p-2} u,  &&\text { in } \Omega, \\
u&=0,
&&\text { on } \partial \Omega,
\end{aligned}
\end{equation}
in the weak (distributional) sense.

\begin{defn}\label{def:eigenfunction}
A nontrivial function $u \in S_{0}^{1,p}(\Omega)$ is called an eigenfunction of problem \eqref{eigenvalue problem} if
\begin{equation}\label{eigenfunction}
\int_{\Omega}|X u|^{p-2} Xu \cdot X \varphi d x=\lambda \int_{\Omega}|u|^{p-2} u \varphi d x
\end{equation}
holds for all $\varphi\in S_{0}^{1,p}(\Omega)$. Such $\lambda$ is called an eigenvalue.
\end{defn}

\begin{thm}\label{positivity of eigenvalues}
All eigenvalues are   positive.
\end{thm}
\begin{proof}
We claim that
\begin{equation*}
\int_{\Omega}|Xu|^{p}dx>0.
\end{equation*}
However, by contradiction, suppose $\int_{\Omega}|Xu|^{p}dx=0$. Then $X_{j}u=0$ a.e. in $\Omega$ for $j=1,...,m$. Since every element of a (Jacobian) basis can be represented by commutators of $\{X_{1},...,X_{m}\}$, it follows that $u=\text{const.}$ a.e. in $\Omega$. Since $u=0$ on $\partial\Omega$, it turns out that $u=0$ a.e. in $\Omega$. By continuity (see Corollary \ref{Holder continuity} which is independent from Theorem \ref{positivity of eigenvalues}) $u=0$ in $\Omega$. Therefore, $u$ cannot be an eigenfunction which leads to a contradiction. 
\end{proof}

Now we estimate a lower bound of eigenvalues which will be used in the proofs of Lemma \ref{L infty boundedness} and Theorem \ref{isolated theorem}.
\begin{lem} \label{lower bound lemma}
Let $(\lambda,u)$ be an eigenpair. Then $\lambda$ has a lower bound
\begin{equation}\label{lower bound of eigenvalues}
\lambda \geq  \frac{1}{C|\Omega|^{1-\frac{p}{p^{*}}}}, 
\end{equation}
where $p^{*}$ is expressed as in \eqref{p asterisk} and the constant $C>0$ does not depend on $u$ and $\lambda$.
\end{lem}
\begin{proof}
Let $(u,\lambda)$ be an eigenpair of problem \eqref{eigenvalue problem}. Using the H\"older inequality and suitable Sobolev embeddings in Theorem \ref{compact Sobolev embedding}
\begin{equation*}
S_{0}^{1,p}(\Omega)\hookrightarrow\begin{cases}
L^{\frac{Qp}{Q-p}}(\Omega) & \text{ if } 1<p<Q,\\
L^{2p}(\Omega) & \text{ if } p\geq Q,
\end{cases}
\end{equation*}
we obtain
\begin{equation*}
\int_{\Omega}|X u|^{p} d x= \lambda\int_{\Omega}|u|^{p} d x\leq \lambda\|u\|_{p^{*}}^{p}|\Omega|^{1-\frac{p}{p^{*}}}\leq C\lambda|\Omega|^{1-\frac{p}{p^{*}}}\int_{\Omega}|X u|^{p} d x ,
\end{equation*}
where 
\begin{equation}\label{p asterisk}
p^{*}=\begin{cases}
\frac{Qp}{Q-p},& \text{ if } 1<p<Q,\\
2p,& \text{ if } p\geq Q.
\end{cases}
\end{equation}

\end{proof}



\begin{thm}\label{closed set}
The set of eigenvalues is closed.
\end{thm}
\begin{proof} 
Let $\{\mu_{j}\}$ be a  sequence of eigenvalues converging to $\mu\neq \infty$. The corresponding  eigenfunctions are denoted by $\{v_{j}\}$. We see that if $v_{j}$ is an eigenfunction, then so is $\frac{1}{\|v_{j}\|_{p}}v_{j}$. Hence, without loss of generality, we assume that
$\left\|v_{j}\right\|_{p}=1$.

By the Eberlein-Šmulian theorem and Theorem \ref{compact Sobolev embedding}, there exist a subsequence $\{v_{j_{k}}\}$ and a function $v\in S_{0}^{1,p}(\Omega)$ such that
\begin{equation}\label{strong and weak convergence}
v_{j_{k}} \to v \text{ strongly in } L^{p}(\Omega)\quad\text{and} \quad X v_{j_{k}} \rightharpoonup  X v \text{ weakly in } L^{p}(\Omega). 
\end{equation}
Regarding $v_{j_{k}}$ and $v$ as test functions, we insert them into \eqref{eigenfunction} to get
\begin{equation*}
\begin{aligned}
&\int_{\Omega}\left(\left|X v_{j_{k}}\right|^{p-2} X v_{j_{k}}-|X v|^{p-2} X v\right) \cdot \left(Xv_{j_{k}}-Xv\right) d x \\
&=\mu_{j_{k}} \int_{\Omega}\left|v_{j_{k}}\right|^{p-2} v_{j_{k}}\left(v_{j_{k}}-v\right) d x-\int_{\Omega}|X v|^{p-2} X v \cdot\left(X v_{j_{k}}-X v\right) d x .
\end{aligned}
\end{equation*}
Taking into account \eqref{strong and weak convergence} and letting $k\to \infty$, the above equality yields 
\begin{equation}\label{closed set: limit}
\lim _{k \rightarrow \infty} \int_{\Omega}\left(\left|X v_{j_{k}}\right|^{p-2} X v_{j_{k}}-|X v|^{p-2} X v\right) \cdot\left(X v_{j_{k}}-X v\right) d x=0 .
\end{equation}
If one proves the following convergence
\begin{equation}\label{closed set: strong convergence of the horizontal gradient}
X v_{j_{k}} \to  X v\quad \text{ strongly in } L^{p}(\Omega), 
\end{equation}
then $v$ is an eigenfunction corresponding to $\mu$ and the proof is complete. 

In order to prove convergence \eqref{closed set: strong convergence of the horizontal gradient}, we treat $1<p<2$ and $p\geq2$ cases separately.

\textbf{Case 1}

Let $p\geq 2$. A modification of inequality \eqref{convexity (a)} yields
\begin{equation}\label{closed set: inequality}
C\left|w_{2}-w_{1}\right|^{p} \leq p\left(\left|w_{2}\right|^{p-2} w_{2}-\left|w_{1}\right|^{p-2} w_{1}\right) \cdot\left(w_{2}-w_{1}\right),
\end{equation}
for all $w_{1}, w_{2} \in \mathbb{R}^{n}$. Inserting $w_{1}=Xv$ and $w_{2}=X v_{j_{k}}$ into \eqref{closed set: inequality}, we see that \eqref{closed set: limit} implies the strong convergence $\{X v_{j_{k}}\}$ in $L^{p}$-norm 
$$
\lim _{k \rightarrow \infty} \int_{\Omega}\left|X v_{j_{k}}-X v\right|^{p} d x=0 .
$$

\textbf{Case 2}

Let $1<p<2$.
A modification of inequality \eqref{convexity (b)} leads to
\begin{equation*}
C\frac{|w_{2}-w_{1}|^{2}}{(|w_{1}|+|w_{2}|)^{2-p}}   \leq p\left(\left|w_{2}\right|^{p-2} w_{2}-\left|w_{1}\right|^{p-2} w_{1}\right)\cdot\left(w_{2}-w_{1}\right) \quad \text{ for all }w_{1}, w_{2} \in \mathbb{R}^{n},
\end{equation*}
so that
\begin{equation*}
\lim_{k\to\infty}\int_{\Omega}\frac{|X v_{j_{k}}-X v|^{2}}{(|X v|+|X v_{j_{k}}|)^{2-p}}dx=0.
\end{equation*}
To get the strong convergence $\{X v_{j_{k}}\}$ in $L^{p}$-norm, we apply the H\"older inequality backward
\begin{equation*}
\begin{aligned}
&\int_{\Omega}|X v_{j_{k}}-X v|^{p}dx=\int_{\Omega}\frac{|X v_{j_{k}}-X v|^{p}}{(|X v|+|X v_{j_{k}}|)^{\frac{p(2-p)}{2}}}(|X v|+|X v_{j_{k}}|\big)^{\frac{p(2-p)}{2}}dx
\\
&\leq\left(\int_{\Omega}\frac{|X v_{j_{k}}-X v|^{2}}{(|X v|+|X v_{j_{k}}|)^{2-p}}dx \right)^{\frac{p}{2}}
\left(\int_{\Omega}\Big(|X v|^{p}+|X v_{j_{k}}|^{p}\Big)dx \right)^{\frac{2-p}{2}}
\\
&\leq C\left(\int_{\Omega}\frac{|X v_{j_{k}}-X v|^{2}}{(|X v|+|X v_{j_{k}}|)^{2-p}}dx \right)^{\frac{p}{2}}.
\end{aligned}
\end{equation*}
Letting $k\to\infty$ we obtain convergence \eqref{closed set: strong convergence of the horizontal gradient} as long as the right-hand side converges to 0. This completes the proof.
\end{proof}

%There are some unbounded functions in $S^{1,p}_{0}(\Omega)$ and for this reason we will find an upper bound of eigenfunctions.
\begin{lem}\label{L infty boundedness}
Let $(u,\lambda)$ be an eigenpair. Then $u\in L^{\infty}(\Omega)$.
\end{lem}
\begin{proof}
For a fixed $k>0$ define a function
\begin{equation*}
\eta(x):=\max_{x\in\Omega} \{u(x), k\},
\end{equation*}
which is in $S^{1,p}_{0}(\Omega)$ and can be regarded as a test function  so that
\begin{equation}\label{upper bound: weak solution}
\int_{\Omega_{k}}|X u|^{p} d x=\lambda \int_{\Omega_{k}}|u|^{p-2} u(u-k) d x    
\end{equation}
holds, where
\begin{equation*}
\Omega_{k}:=\{x \in \Omega \mid u(x)>k\} .
\end{equation*}
Also, observe that 
\begin{equation}\label{upper bound: some properties}
k\left|\Omega_{k}\right| \leq\|u\|_{1}\quad  \text{and }\quad  \lim_{k\to \infty}\left|\Omega_{k}\right|=0.
\end{equation}
First, by using a consequence of  the Clarkson inequality 
\begin{equation*}
a^{p-1} \leq 2^{p-1}(a-k)^{p-1}+2^{p-1} k^{p-1}, \quad 1<p<\infty, 
\end{equation*}
to the right-hand side of \eqref{upper bound: weak solution},  we get
\begin{equation}\label{upper bound: Clarkson's inequality}
\int_{\Omega_{k}} u^{p-1}(u-k) d x \leq 2^{p-1} \int_{\Omega_{k}}(u-k)^{p} d x+2^{p-1} k^{p-1} \int_{\Omega_{k}}(u-k) d x.
\end{equation}
Then we apply Poincar\'e inequality \eqref{Poincare inequality}  and lower bound \eqref{lower bound of eigenvalues} to obtain
\begin{equation}\label{upper bound: Poincare and lower bound}
\begin{aligned}
\int_{\Omega_{k}}(u-k)^{p} d x \leq\lambda_{1}^{-1} \int_{\Omega_{k}}|X u|^{p} d x \leq C \left|\Omega_{k}\right|^{1-\frac{p}{p^{*}}} \int_{\Omega_{k}}|X u|^{p} d x\\
=\lambda  C \left|\Omega_{k}\right|^{1-\frac{p}{p^{*}}}\int_{\Omega_{k}}u^{p-1} (u-k) d x,
\end{aligned}
\end{equation}
where $p^{*}$ is expressed as in \eqref{p asterisk}. Multiplying both sides of \eqref{upper bound: Clarkson's inequality} by $\lambda  C \left|\Omega_{k}\right|^{1-\frac{p}{p^{*}}}$ and using inequality \eqref{upper bound: Poincare and lower bound} we obtain
\begin{equation*}
\left[1-2^{p-1}C \lambda\left|\Omega_{k}\right|^{1-\frac{p}{p^{*}}}\right] \int_{\Omega_{k}}(u-k)^{p} d x \leq 2^{p-1}Ck^{p-1} \lambda\left|\Omega_{k}\right|^{1-\frac{p}{p^{*}}} \int_{\Omega_{k}}(u-k) d x .
\end{equation*}
By virtue of \eqref{upper bound: some properties}, we find a bound
\begin{equation*}
2^{p-1}C \lambda\left|\Omega_{k}\right|^{1-\frac{p}{p^{*}}} \leq \frac{1}{2}, \quad \text{ for all }    k \geq k_{1}=2^{\frac{p^{*}(p-2)}{ p^{*}-p}}C^{{\frac{p^{*}}{ p^{*}-p}}} \lambda^{{\frac{p^{*}}{ p^{*}-p}}}\|u\|_{1}.
\end{equation*}
Therefore,
\begin{equation}\label{upper bound: Lp-L1 inequality}
\int_{\Omega_{k}}(u-k)^{p} d x \leq 2^{p}C k^{p-1} \lambda\left|\Omega_{k}\right|^{1-\frac{p}{p^{*}}} \int_{\Omega_{k}}(u-k) d x  \quad \text{ for all } k \geq k_{1}.
\end{equation}
Finally, applying the H\"older inequality backward to the left-hand side of \eqref{upper bound: Lp-L1 inequality}, we arrive at
\begin{equation*}
\int_{\Omega_{k}}(u-k) d x \leq  (2^{p}C\lambda)^{\frac{1}{p-1}} k\left|\Omega_{k}\right|^{1+\frac{p^{*}-p}{p^{*}(p-1)}} \quad \text{ for all } k \geq k_{1}.  
\end{equation*}
This is the desired inequality to bound $\esssup_{\Omega}u$ (cf \cite{Lin1992}).
\end{proof}

\begin{cor}\label{Holder continuity}
All eigenfunctions  are H\"older continuous.
\end{cor}
\begin{proof}
Let $u$ be an eigenfunction. The proof is split in two cases.

Let $1<p\leq Q$. Since $u$ is bounded in $L^{\infty}(\Omega)$ by Lemma \ref{L infty boundedness}, the proof follows from \cite[Theorem 6.6]{DGN1998}.

Let $p>Q$. Then $u$ is  H\"older continuous with exponent $\alpha=1-\frac{Q}{p}$ with respect to the control distance $d_{X}$  by \cite[Theorem 1.11]{GN1996}.
\end{proof}

\begin{thm}\label{existence and uniqueness of the Rayleigh quotient}
There exists a unique nontrivial normalized $u\in S_{0}^{1,p}(\Omega)$  such that
 $$\|Xu\|_{p}^{p}=\inf_{\stackrel{v \in S_{0}^{1,p}(\Omega)}{\|v\|_{p}=1}}\|Xv\|_{p}^{p}.$$
\end{thm}
\begin{proof}
The proof is standard, i.e., relies on weakly lower semicontinuity of \eqref{Poincare norm}, Poincar\'e inequality \eqref{Poincare inequality} and reflexivity of $S_{0}^{1,p} (\Omega)$.
\end{proof} 

By virtue of Theorem \ref{existence and uniqueness of the Rayleigh quotient}, we obtain the smallest eigenvalue.
\begin{thm}\label{minimization theorem}
There exists a nontrivial function $u_{1}\in S_{0}^{1,p}(\Omega)$ such that 
\begin{equation*}
\inf_{v \in S_{0}^{1,p}(\Omega)\setminus\{0\}}\left\{\frac{\int_{\Omega}|X v|^{p} d x}{\int_{\Omega}|v|^{p} d x}\right\}=\lambda_{1}=\frac{\int_{\Omega}|X u_{1}|^{p} d x}{\int_{\Omega}|u_{1}|^{p} d x}.
\end{equation*}
Here $(\lambda_{1}, u_{1})$ is the first eigenpair.
\end{thm}
\begin{proof}
Consider a functional defined by
\begin{equation*}
f(\varepsilon):=\frac{\int_{\Omega}|X(u_{1}+\varepsilon v)|^{p} d x}{\int_{\Omega}|u_{1}+\varepsilon v|^{p} d x},
\end{equation*}
where $v\in S_{0}^{1,p}(\Omega)$ and  $u_{1}$ minimizes the Rayleigh quotient by Theorem \ref{existence and uniqueness of the Rayleigh quotient}. Then $f$ must satisfy $f^{\prime}(0)=0$. 
We see that the expression $f^{\prime}(0)=0$ is equivalent to
\begin{equation*}
\int_{\Omega}|u_{1}|^{p} d x \int_{\Omega}|X u_{1}|^{p-2} X u_{1} \cdot X v d x=\int_{\Omega}|u_{1}|^{p-2} u_{1} v d x \int_{\Omega}|X u_{1}|^{p} d x,
\end{equation*}
which implies
\begin{equation*}
\begin{aligned}
\int_{\Omega}|X u_{1}|^{p-2} X u_{1} \cdot X v d x &=\frac{\int_{\Omega}|X u_{1}|^{p} d x}{\int_{\Omega}|u_{1}|^{p} d x} \int_{\Omega}|u_{1}|^{p-2} u_{1} v d x\\
&=\lambda_{1} \int_{\Omega}|u_{1}|^{p-2} u_{1} v d x.
\end{aligned}
\end{equation*}
So, $(\lambda_{1}, u_{1})$ is an eigenpair. Finally, by the contradiction argument it can be shown that $\lambda_{1}$ is the smallest.
\end{proof}
\begin{rem}\label{the best constant of Poincare}
Theorem \ref{minimization theorem} also yields the best constant $C=\lambda_{1}^{-1}$ of  Poincar\'e inequality \eqref{Poincare inequality}.
\end{rem}

As a consequence of Theorem \ref{minimization theorem}, we prove the nonnegativity of $u_{1}$.
\begin{cor}
Let $u_{1}$ be a minimizer in Theorem \ref{minimization theorem}. Then $|u_{1}|$ also minimizes the Rayleigh quotient.
\end{cor}
\begin{proof}
Proposition \ref{proposition Xu} together with Corollary \ref{Holder continuity} ensure that $|Xu_{1}|=|X|u_{1}||$  a.e. in $\Omega$.
\end{proof}

\begin{cor}\label{Harnack inequality}
Let $u_{1}$ be a minimizer in Theorem \ref{minimization theorem}. Then $|u_{1}|>0$ in $\Omega$.
\end{cor}
\begin{proof}
Harnack inequality \cite[Theorem 3.9]{Lu1996} together with Corollary \ref{Holder continuity}
 assert that 
$$|u_{1}|>0 \quad \text{ in } \Omega$$ 
or 
$$|u_{1}|\equiv 0 \quad \text{  in } \Omega.$$  
The latter  contradicts Definition \ref{def:eigenfunction}, so  $|u_{1}|>0$ in $\Omega$. 
\end{proof}


\section{Simplicity and isolatedness of the first eigenvalue}\label{sec4}
\begin{thm}\label{simplicity theorem}
The first eigenvalue $\lambda_{1}$ is simple. In other words, if there exist two positive eigenfunctions $u$ and $v$ associated with $\lambda_{1}$, then they are proportional.
\end{thm}
\begin{proof}
The proof is an adaptation of the $p$-Laplacian case \cite{Lin1990}. Suppose there exist two positive first eigenfunctions $u$ and $v$ corresponding to $\lambda_{1}$.
Set two test functions
\begin{equation*}
\varphi=\frac{(u+\varepsilon)^{p}-(v+\varepsilon)^{p}}{(u+\varepsilon)^{p-1}} \text { for $u$ and } \psi=\frac{(v+\varepsilon)^{p}-(u+\varepsilon)^{p}}{(v+\varepsilon)^{p-1}} \text{ for $v$,}
\end{equation*}
where $\varepsilon>0$ is arbitrary. 
For convenience, we also set
\begin{equation*}
u_{\varepsilon}=u+\varepsilon, \quad v_{\varepsilon}=v+\varepsilon  .  
\end{equation*}
Substituting $u,\varphi$ and $v,\psi$ into equations \eqref{eigenfunction} respectively, and then adding these equations we compute
\begin{equation*}
\begin{aligned}
&\lambda_{1} \int_{\Omega}  {\left[\frac{u^{p-1}}{u_{\varepsilon}^{p-1}}-\frac{v^{p-1}}{v_{\varepsilon}^{p-1}}\right]\left(u_{\varepsilon}^{p}-v_{\varepsilon}^{p}\right) d x } \\
&= \int_{\Omega}\left[\left\{1+(p-1)\left(\frac{v_{\varepsilon}}{u_{\varepsilon}}\right)^{p}\right\}\left|X u_{\varepsilon}\right|^{p}+\left\{1+(p-1)\left(\frac{u_{\varepsilon}}{v_{\varepsilon}}\right)^{p}\right\}\left|X v_{\varepsilon}\right|^{p}\right] d x \\
&-\int_{\Omega}\left[p\left(\frac{v_{\varepsilon}}{u_{\varepsilon}}\right)^{p-1}\left|X u_{\varepsilon}\right|^{p-2} X u_{\varepsilon} \cdot X v_{\varepsilon}+p\left(\frac{u_{\varepsilon}}{v_{\varepsilon}}\right)^{p-1}\left|X v_{\varepsilon}\right|^{p-2} X v_{\varepsilon} \cdot X u_{\varepsilon}\right] d x \\
&= \int_{\Omega}\left(u_{\varepsilon}^{p}-v_{\varepsilon}^{p}\right)\left(\left|X \log u_{\varepsilon}\right|^{p}-\left|X \log v_{\varepsilon}\right|^{p}\right) d x \\
&-\int_{\Omega} p v_{\varepsilon}^{p}\left|X \log u_{\varepsilon}\right|^{p-2} X \log u_{\varepsilon} \cdot\left(X \log v_{\varepsilon}-X \log u_{\varepsilon}\right) d x \\
&-\int_{\Omega} p u_{\varepsilon}^{p}\left|X \log v_{\varepsilon}\right|^{p-2} X \log v_{\varepsilon} \cdot\left(X \log u_{\varepsilon}-X \log v_{\varepsilon}\right) d x.
\end{aligned}
\end{equation*}

Now consider the case $p\geq 2$.
According to inequality  \eqref{convexity (a)} with $w_{1}=X \log v_{\varepsilon}$ and $w_{2}=X \log u_{\varepsilon}$, we have
\begin{equation}\label{simplicity: convexity (a) vu}
\begin{aligned}
0  &\leq C \int_{\Omega}\frac{1}{v_{\varepsilon}^{p}}\left|v_{\varepsilon} X u_{\varepsilon}-u_{\varepsilon} X v_{\varepsilon}\right|^{p} d x \\
& \leq \int_{\Omega}u_{\varepsilon}^{p}\left(\left|X \log u_{\varepsilon}\right|^{p}-\left|X \log v_{\varepsilon}\right|^{p}\right) d x\\
&-\int_{\Omega} p u_{\varepsilon}^{p}\left|X \log v_{\varepsilon}\right|^{p-2} X \log v_{\varepsilon} \cdot\left(X \log u_{\varepsilon}-X \log v_{\varepsilon}\right) d x.
\end{aligned}
\end{equation}
According to inequality \eqref{convexity (a)} with $w_{1}=X \log u_{\varepsilon}$ and $w_{2}=X \log v_{\varepsilon}$, we have
\begin{equation}\label{simplicity: convexity (a) uv}
\begin{aligned}
0 & \leq C \int_{\Omega}\frac{1}{u_{\varepsilon}^{p}}\left|v_{\varepsilon} X u_{\varepsilon}-u_{\varepsilon} X v_{\varepsilon}\right|^{p} d x \\
& \leq-\int_{\Omega}v_{\varepsilon}^{p}\left(\left|X \log u_{\varepsilon}\right|^{p}-\left|X \log v_{\varepsilon}\right|^{p}\right) d x\\
& -\int_{\Omega} p v_{\varepsilon}^{p}\left|X \log u_{\varepsilon}\right|^{p-2} X \log u_{\varepsilon} \cdot\left(X \log v_{\varepsilon}-X \log u_{\varepsilon}\right) d x.
\end{aligned}   
\end{equation}
Combining \eqref{simplicity: convexity (a) vu} and \eqref{simplicity: convexity (a) uv}, we obtain
\begin{equation*}
\begin{aligned}
&0  \leq C \int_{\Omega}\left(\frac{1}{v_{\varepsilon}^{p}}+\frac{1}{u_{\varepsilon}^{p}}\right)\left|v_{\varepsilon} X u_{\varepsilon}-u_{\varepsilon} X v_{\varepsilon}\right|^{p} d x \\
& \leq\lambda_{1} \int_{\Omega}\left[\left(\frac{u}{u_{\varepsilon}}\right)^{p-1}-\left(\frac{v}{v_{\varepsilon}}\right)^{p-1}\right]\left(u_{\varepsilon}^{p}-v_{\varepsilon}^{p}\right) d x.
\end{aligned}    
\end{equation*}
Observe that by the dominated convergence theorem we have
\begin{equation}\label{simplicity: dominated convergence}
\lim _{\varepsilon \rightarrow 0+} \int_{\Omega}\left[\frac{u^{p-1}}{u_{\varepsilon}^{p-1}}-\frac{v^{p-1}}{v_{\varepsilon}^{p-1}}\right]\left(u_{\varepsilon}^{p}-v_{\varepsilon}^{p}\right) d x=0    .
\end{equation}
Taking into account \eqref{simplicity: dominated convergence}, we apply the Fatou lemma as $\varepsilon\to 0$ to obtain
\begin{equation*}
\int_{\Omega}\left(\frac{1}{v^{p}}+\frac{1}{u^{p}}\right)\left|v X u-u X v\right|^{p} d x=0.    
\end{equation*}
This is equivalent to  $X\left(\dfrac{u}{v}\right)=0$ a.e. in $\Omega$.  Therefore, by continuity (Corollary \ref{Holder continuity}) $u=cv$ for some $c\neq 0$. 

The proof of case $1<p<2$ is similar. The only difference is that we use inequality \eqref{convexity (b)} twice by interchanging $\omega_{1}$ and $\omega_{2}$ to obtain

\begin{equation*}
\begin{aligned}
0 & \leq C \int_{\Omega}\left(u_{\varepsilon} v_{\varepsilon}\right)^{p}\left(u_{\varepsilon}^{p}+v_{\varepsilon}^{p}\right) \frac{\left|v_{\varepsilon} X u_{\varepsilon}-u_{\varepsilon} X v_{\varepsilon}\right|^{2}}{\left(v_{\varepsilon}\left|X u_{\varepsilon}\right|+u_{\varepsilon}\left|X v_{\varepsilon}\right|\right)^{2-p}} d x \\
& \leq\lambda_{1} \int_{\Omega}\left[\left(\frac{u}{u_{\varepsilon}}\right)^{p-1}-\left(\frac{v}{v_{\varepsilon}}\right)^{p-1}\right]\left(u_{\varepsilon}^{p}-v_{\varepsilon}^{p}\right) d x.
\end{aligned}    
\end{equation*}
Again, applying the dominated convergence theorem and Fatou lemma as $\varepsilon\to 0$ we arrive at $X\left(\dfrac{u}{v}\right)=0$ a.e. in $\Omega$. Then $u=cv$ for some $c\neq 0$.
\end{proof}

The techniques used in Theorem \ref{simplicity theorem} allow us to prove Theorem \ref{change sign theorem} and Theorem \ref{isolated theorem} (cf \cite{AL1987}).
\begin{thm}\label{change sign theorem}
For every eigenvalue $\lambda\neq \lambda_{1}$  all corresponding eigenfunctions change their sign in $\Omega$.
\end{thm}
\begin{proof}
Let $v$ be any eigenfunction corresponding to $\lambda\neq\lambda_{1}$, i.e., $\lambda> \lambda_{1}$. Also for convenience, in the proof we denote the first eigenfunction $u_{1}$ by $u$. Suppose $v\geq 0$. Proceeding the same procedure as in the proof of Theorem \ref{simplicity theorem} we obtain
\begin{equation*}
\int_{\Omega}\left[\lambda_{1}\left(\frac{u}{u_{\varepsilon}}\right)^{p-1}-\lambda\left(\frac{v}{v_{\varepsilon}}\right)^{p-1}\right]\left(u_{\varepsilon}^{p}-v_{\varepsilon}^{p}\right) d x \geq 0.    
\end{equation*}
Then
\begin{equation*}
\int_{\Omega}\left(u^{p}-v^{p}\right) d x \leq 0,  
\end{equation*}
as $\varepsilon\to 0$. One can multiply $u$ by a sufficiently large positive constant by Theorem \ref{simplicity theorem} to derive $u^{p}-v^{p}>0$ in $\Omega$. However, this leads to a contradiction. Thus, $v$ changes sign in $\Omega$.
\end{proof}

\begin{thm}\label{isolated theorem}
The first eigenvalue $\lambda_{1}$ is isolated.
\end{thm}
\begin{proof}
Let $\mu\neq \lambda_{1}$ be any eigenvalue of problem \eqref{eigenvalue problem}. Then $\mu > \lambda_{1}$ by Theorem \ref{minimization theorem}. Let $\{v_{j}\}$ be  eigenfunctions associated with $\mu_{j}$ such that $\|v_{j}\|_{p}=1$. Thus, we get
\begin{equation*}
\mu_{j}=\int_{\Omega}|Xv_{j}|^{p}dx. 
\end{equation*}
By the Eberlein-Šmulian theorem and Theorem \ref{compact Sobolev embedding}, there exist a subsequence $\{v_{j_{k}}\}$ and a function $v\in S_{0}^{1,p}(\Omega)$ such that
\begin{equation*}
v_{j_{k}} \to v\enskip \text{ strongly in } L^{p}(\Omega)\quad \text{and} \quad X v_{j_{k}} \rightharpoonup  X v \enskip \text{ weakly in } L^{p}(\Omega). 
\end{equation*}
In view of the weakly lower semicontinuity of $\|\cdot\|_{p}$, we have
\begin{equation*}
\int_{\Omega}|Xv|^{p}dx \leq \lim_{k\to\infty}\mu_{j_{k}} =\lambda_{1}, 
\end{equation*}
so $v$ minimizes the Rayleigh quotient and is the first eigenfunction by Theorem \ref{minimization theorem} . Then $v>0$ in $\Omega$ by Corollary \ref{Harnack inequality} .

By contradiction suppose there exists $N_{1}\in\mathbb{N}$ such that $\mu_{j_{k}} \neq \lambda_{1}$ for all $j_{k}\geq N_{1}$, then $v_{j_{k}}$ changes signs in $\Omega$ by Theorem \ref{change sign theorem}. Therefore, $\Omega_{j_{k}}^{+}:=\left\{v_{j_{k}}>0\right\}$ and $\Omega_{j_{k}}^{-}:=\left\{v_{j_{k}}<0\right\}$ are nonempty sets.

The lower bound \eqref{lower bound of eigenvalues} gives
$$
\left|\Omega_{j_{k}}^{+}\right| \geq \left(C\mu_{j_{k}}\right)^{\frac{p^{*}}{p^{*}-p}}\quad \text{and} \quad\left|\Omega_{j_{k}}^{-}\right| \geq \left(C\mu_{j_{k}}\right)^{\frac{p^{*}}{p^{*}-p}},$$
where $p^{*}$ is expressed as in \eqref{p asterisk}.  Consequently, we obtain
$$\left|\Omega^{\pm}\right|=\left|\limsup\Omega_{j_{k}}^{\pm}\right|> 0.$$ 
By \cite[Theorem 4.9]{Bre2011}  there exists $\{v_{j_{k_{l}}}\}\subset \{v_{j_{k}}\}$ such that $v_{j_{k_{l}}}\to v$ a.e. in $\Omega$. It turns out that $v>0$ a.e. in $\Omega^{+}$ and $v<0$ a.e. in $\Omega^{-}$ as $j_{k_{l}}\to \infty$, i.e., the first eigenfunction changes sign in $\Omega$. This is a contradiction, so $\lambda_{1}$ is isolated.
\end{proof}

\section{Appendix}
The positive and negative parts of a function $u$ are given by
$$u^{+}(x)=\max \{u(x), 0\}\quad\text{and}\quad u^{-}(x)=-\min \{u(x), 0\}.$$ Also they satisfy $u=u^{+}-u^{-}$ and $|u|=u^{+}+u^{-}$.
\begin{prop}\label{proposition Xu} \cite[Lemma 3.5]{GN1996}. 
Let $F\in C^{1}({\mathbb{R}})$ such that $F'\in L^{\infty}({\mathbb{R}})$. Then for all $u \in S_{0}^{1,p}(\Omega)$  we have 
\begin{enumerate}[label=(\roman*)]
    \item The composite function $h=F\circ u$ belongs to $S_{0}^{1,p}(\Omega)$ and for all $\varphi\in C^{\infty}_{c}(\Omega)$ and for $j=1,...,m$ we have
    \begin{equation*}
       \int_{\Omega} X_{j}(F\circ u)\varphi dx=\int_{\Omega}(F\circ u)X_{j}u\varphi dx.
    \end{equation*}
    \item Moreover, $u^{\pm}, |u| \in S_{0}^{1,p}(\Omega)$ and 
\begin{equation*}
Xu^{+}=\begin{cases}
Xu, & \text{if $u>0$ a.e. in $\Omega$,}\\
0, & \text{if $u\leq 0$ a.e.  in $\Omega$.}
\end{cases}
\end{equation*}
\begin{equation*}
Xu^{-}=\begin{cases}
0, & \text{if $u\geq0$ a.e.  in $\Omega$,}\\
-Xu & \text{if $u<0$ a.e.  in $\Omega$.}
\end{cases}
\end{equation*}
\begin{equation*}
X|u|=\begin{cases}
Xu, & \text{if $u>0$ a.e.  in $\Omega$,}\\
0, & \text{if $u= 0$ a.e.  in $\Omega$,}\\
-Xu & \text{if $u<0$ a.e.  in $\Omega$.}
\end{cases}
\end{equation*}
\end{enumerate}
\end{prop}

\begin{prop}\label{convexity inequality}\cite{Lin1990}
\begin{enumerate}[label=(\roman*)]
\item  Let $p \geq 2$. Then there exists a constant $C>0$ such that
\begin{equation}\label{convexity (a)}
|w_{2}|^{p} \geq|w_{1}|^{p}+p|w_{1}|^{p-2} w_{1} \cdot(w_{2}-w_{1})+C|w_{1}-w_{2}|^{p}
\end{equation}
holds for all $w_{1}, w_{2} \in \mathbb{R}^{n}$. The constant $C$ depends on $p$.

\item  Let $1<p<2$. Then there exists a constant $C>0$ such that
\begin{equation}\label{convexity (b)}
|w_{2}|^{p} \geq|w_{1}|^{p}+p|w_{1}|^{p-2} w_{1} \cdot(w_{2}-w_{1})+C \frac{|w_{1}-w_{2}|^{2}}{(|w_{1}|+|w_{2}|)^{2-p}} 
\end{equation}
holds  for all $w_{1}, w_{2} \in \mathbb{R}^{n}$. The constant $C$ depends on $p$.
\end{enumerate}
\end{prop}

\section{Acknowledgments}

This research is funded by the Science Committee of the Ministry of Science and Higher Education of Kazakhstan (Grant No. AP09057887). This work was partially supported by the NU grant 20122022FD4105. No new data was collected or generated during the course of research.



\begin{thebibliography}{HOHOLT08}




\bibitem{AL1987}
A. Anane  and J. L. Lions. Simplicité et isolation de la premiere valeur propre du p-laplacien avec poids. {\em Comptes rendus de l'Académie des sciences. Série 1, Mathématique}, 305(16), pp. 725-728, 1987.

\bibitem{Bre2011}
H. Brézis. Functional analysis, Sobolev spaces and partial differential equations.  Springer New York, 2011.

\bibitem{CG2022}
M. Carfagnini and M. Gordina. Dirichlet sub-Laplacians on homogeneous Carnot groups: spectral properties, asymptotics, and heat content. {\em International Mathematics Research Notices}, https://doi.org/10.1093/imrn/rnad065, 2023.

\bibitem{Chen2021}
H. Chen and H.G. Chen. Estimates of Dirichlet eigenvalues for a class of sub‐elliptic operators. {\em Proceedings of the London Mathematical Society}, 122(6), pp.808-847, 2021.

\bibitem{DGN1998}
D. Danielli, N. Garofalo and D. M. Nhieu. Trace inequalities for Carnot-Carathéodory spaces and applications. {\em Annali della Scuola Normale Superiore di Pisa-Classe di Scienze}, 27(2), pp. 195-252, 1998.


\bibitem{GN1996}
N. Garofalo and D. M. Nhieu.  Isoperimetric and Sobolev inequalities for Carnot‐Carathéodory spaces and the existence of minimal surfaces. {\em Communications on Pure and Applied Mathematics}, 49(10), pp. 1081-1144, 1996.

\bibitem{Jer1986}
D. Jerison. The Poincaré inequality for vector fields satisfying Hörmander’s condition. {\em Duke Mathematical Journal}, 53(2), pp. 503-523, 1986.

\bibitem{Lin1990}
P. Lindqvist.
On the Equation $\operatorname{div}(|\nabla u|^{p-2}\nabla u) + \lambda|u| ^{p-2}u = 0$.
{\em Proceedings of the American Mathematical Society}, Vol. 109, No. 1, pp. 157-164, 1990.

\bibitem{Lin1992}
P. Lindqvist.
Addendum: “On the equation $\textrm {div}(\vert \nabla u\vert ^ {p-2}\nabla u)+\lambda \vert u\vert ^ {p-2}u=0$”. {\em Proceedings of the American Mathematical Society}, Vol. 116, No. 2, pp. 583-584. 

\bibitem{Lu1996}
G. Lu. Embedding theorems into Lipschitz and BMO spaces and applications to quasilinear subelliptic differential equations. {\em Publicacions Matemátiques}, pp. 301-329, 1996

\bibitem{NSW1985}
A. Nagel, E.M. Stein and S. Wainger. Balls and metrics defined by vector fields I: Basic properties. {\em Acta Mathematica}, 155, pp. 103-147, 1985





\bibitem{Str1986}
R.S. Strichartz. Sub-Riemannian geometry. {\em Journal of Differential Geometry}, 24(2), pp. 221-263, 1986.


\bibitem{WNL2009}
N. Wei, P. Niu and H. Liu. Dirichlet Eigenvalue Ratios for the $p$-sub-Laplacian in the Carnot Group. {\em Journal of Partial Differential Equations}, 22(1), pp. 1-10, 2009.
\end{thebibliography}
\end{document}

